\documentclass[../cours_2de.tex]{subfiles}
\begin{document}
\setcounter{chapter}{17}
\chapter{Statistiques : position et lecture critique}
\label{sec:ch18}
\begin{objectif}
Décrire une série ; interpréter moyenne, médiane et quartiles ; lire des graphiques et mesurer l'effet d'une donnée ajoutée ou supprimée.
\end{objectif}
\section{Population, variable et distribution}
Une population est l'ensemble étudié, un individu une unité statistique.
Une variable qualitative décrit une catégorie ; une variable quantitative une quantité, discrète ou continue.
L'effectif $n_i$ d'une valeur $x_i$ est son nombre d'apparitions ; sa fréquence est $f_i=n_i/N$ avec $N=\sum n_i>0$.
Ainsi $\sum f_i=1$.
Une étude décrit la population observée ; extrapoler à une population plus grande demande un protocole d'échantillonnage adapté.
\section{Moyenne et linéarité}
\begin{definition}
La moyenne d'une série de $N$ valeurs est
$\overline x=(x_1+\cdots+x_N)/N$.
Pour des valeurs regroupées avec effectifs $n_i$,
$\overline x=\sum n_ix_i/N=\sum f_ix_i$.
\end{definition}
La formule pondérée ne change pas le calcul : elle regroupe les valeurs répétées.
\begin{propriete}
Si chaque valeur est transformée en $y_i=ax_i+b$, alors
$\overline y=a\overline x+b$.
\end{propriete}
\begin{demonstration}
$\sum_{i=1}^N(ax_i+b)=a\sum_{i=1}^Nx_i+Nb$ ; diviser par $N$ donne le résultat.
\end{demonstration}
Si deux groupes ont des effectifs $N_1,N_2$ et des moyennes $m_1,m_2$, leur moyenne globale est
$(N_1m_1+N_2m_2)/(N_1+N_2)$, puisque $N_jm_j$ reconstitue la somme du groupe.
\begin{exemple}
Sur la série ordonnée $1,2,2,4,5,8,8,10$, la somme vaut $40$, donc la moyenne vaut $5$.
Ajouter $3$ à chaque valeur donne une moyenne de $8$.
\end{exemple}
\section{Médiane, quartiles, écart interquartile}
Pour la série ordonnée $x_{(1)}\leq\cdots\leq x_{(N)}$, nous retenons la convention :
si $N$ est impair, médiane de rang $(N+1)/2$ ; si $N$ est pair, moyenne des valeurs de rangs $N/2$ et $N/2+1$.
Une médiane partage la série en deux parties contenant chacune au moins la moitié des observations, égalités comprises.
$Q_1$ est la valeur de rang $\lceil N/4\rceil$ et $Q_3$ celle de rang $\lceil3N/4\rceil$.
L'écart interquartile est $Q_3-Q_1$, l'étendue est maximum moins minimum.
\begin{exemple}
Pour $1,2,2,4,5,8,8,10$, la médiane vaut $(4+5)/2=4{,}5$,
$Q_1=2$, $Q_3=8$, l'écart interquartile vaut $6$ et l'étendue $9$.
\end{exemple}
Ces indicateurs sont des définitions assorties de conventions, pas des théorèmes.
Les logiciels n'utilisent pas toujours la même interpolation pour les quartiles ; annoncer la convention avant une comparaison.
\section{Boîte à moustaches et graphiques}
Une boîte à moustaches représente ici les cinq valeurs minimum, $Q_1$, médiane, $Q_3$, maximum.
\begin{center}
\begin{tikzpicture}[x=.9cm,y=1cm]
\draw[->](0,0)--(11,0);
\foreach \x in {0,1,...,10}{\draw(\x,.05)--(\x,-.05)node[below]{\small\x};}
\draw(1,.7)--(2,.7);\draw(8,.7)--(10,.7);
\draw(1,.45)--(1,.95);\draw(10,.45)--(10,.95);
\draw[fill=accent!10](2,.4)rectangle(8,1);
\draw[thick](4.5,.4)--(4.5,1);
\node[above]at(4.5,1){médiane};
\end{tikzpicture}
\end{center}
Certaines boîtes utilisent des moustaches calculées à partir de $1{,}5$ fois l'écart interquartile : ce n'est pas la convention de cette figure.
Un diagramme en barres compare des catégories ; un secteur circulaire de fréquence $f$ a un angle $360f$ degrés (ou $180f$ pour un demi-disque) par proportionnalité des aires des secteurs.
Une courbe montre une évolution ; un nuage de points montre des couples de mesures, sans prouver de causalité.
\begin{attention}
Lire les axes, les unités, l'origine, la période et la population.
Une échelle verticale tronquée peut amplifier visuellement une différence.
Une corrélation observée ne prouve ni causalité ni absence de variable cachée.
\end{attention}
\section{Ajouter ou supprimer une valeur}
\begin{propriete}
Si une série de taille $N$ et moyenne $m$ reçoit une nouvelle valeur $a$, sa moyenne devient
$m'=(Nm+a)/(N+1)$, donc $m'-m=(a-m)/(N+1)$.
\end{propriete}
\begin{demonstration}
L'ancienne somme vaut $Nm$ ; la nouvelle vaut $Nm+a$ et l'effectif augmente de $1$. En soustrayant $m$ à la nouvelle moyenne, on obtient la différence annoncée.
\end{demonstration}
L'ajout fait monter la moyenne si $a>m$, baisser si $a<m$ ; il la conserve si $a=m$.
Supprimer $a$ d'une série de taille $N>1$ donne $m'=(Nm-a)/(N-1)$.
Pour la médiane, il faut regarder les nouveaux rangs : il n'existe pas de règle analogue fondée sur la seule ancienne médiane.
\begin{exemple}
Ajouter $100$ à la série de l'exemple donne $140/9\simeq15{,}56$ de moyenne, mais une médiane de $5$.
Une valeur extrême peut donc déplacer fortement la moyenne tout en modifiant peu la médiane. Cela ne signifie pas que la médiane est toujours inchangée.
\end{exemple}
\begin{exemple}[Retirer une observation]
Si l'on retire $10$ de la série initiale, il reste $1,2,2,4,5,8,8$ : la moyenne devient $30/7$ et la médiane $4$. Si l'on retire plutôt $1$, la moyenne devient $39/7$ et la médiane $5$. Pour la médiane, le déplacement des rangs doit être étudié après chaque retrait ; une formule utilisant seulement l'ancienne médiane ne suffit pas.
\end{exemple}
\utilite{La moyenne mesure un niveau global ; la médiane résiste mieux aux valeurs extrêmes. Un indicateur ne remplace ni le graphique ni l'examen de la collecte des données.}
\begin{retenir}
Trier avant de chercher médiane et quartiles ; pondérer par les effectifs ; annoncer les conventions et comparer des populations effectivement comparables.
\end{retenir}
\end{document}
